\documentclass[12pt,a4paper]{article}
\usepackage{titlesec}
\usepackage[utf8]{inputenc}
\usepackage[T1]{fontenc}
\usepackage{lmodern}
\usepackage[a4paper]{geometry}
\usepackage[frenchb]{babel}
\usepackage{amsmath, amssymb}
\usepackage{tcolorbox}
\usepackage{framed}
\usepackage{fancybox}
\usepackage{enumitem}
\usepackage{tikz}
\usepackage{setspace}
\onehalfspacing
\pagestyle{plain}
\usepackage{multicol}
\begin{document}
\vspace*{-4cm}
\begin{framed}
\centering \underline{\textbf{Épreuve des MATHÉMATIQUES}} \\
\underline{\textbf{Classe :}} $6^{ième}$ \hfill \underline{\textbf{Durée :}} 1h30min
\end{framed}
\underline{\textbf{Texte :}} Fabrication de savons \\
BAF est une société de fabrication de savons sous deux formes représentées par les figures 1 et 2. Elle est située dans l’angle d’un carrefour de quatre voies $(D1)$, $(D2)$, $(D3)$ et $(D4)$ avec des points de ventes tels qu’indiqué sur le schéma. \\
\includegraphics[width=0.6\textwidth]{ice} \includegraphics[width=0.45\textwidth]{probleme} \\
Marie, une élève en classe de 6ième en visite dans la société se pose plusieurs questions. \\
\underline{\textbf{Tâche :}} Pour ton évaluation, tu es invitée à aider Marie en résolvant les trois problèmes suivants. \\ 
\underline{\textbf{Problème 1}}
\begin{enumerate}
\item Complète le texte suivant en remplaçant les lettres par les mots qui convient : volume ; longueur ; latérales ; l'aire de base du pavé droit ; bases ; carrées ; rectangulaires (utilise uniquement les lettres)\\ \og Un cube est un solide de l'espace dont les six faces sont des (a) et dont les arêtes sont de même (b). Un pavé droit est un solide de l'espace qui a six faces de formes (c), dont quatre faces (d) et deux (e). Pour un pavé droit de longueur $L$, de largeur $l$ et de hauteur $h$, alors l'opération $L \times l$ est la formule permettant de calculer (f) et l'opération $L \times l \times h$ est la formule du (g) d'un pavé droit.\fg
\item 
\begin{enumerate}
\item Donne le nom de chacun des solides représentés par les figures 1 et 2. 
\item Donne la différence qui existe entre ces deux solides représentés par les figues 1 et 2. 
\end{enumerate}
\item Détermine le volume $V$ du solide représenté par la figure 1.
\item Justifie que l'aire totale du solide représenté par la figure 2 est $A_{T}=68cm^{2}$.
\item Représente soigneusement le patron du solide représenté par la figure 2. \\ \\
\underline{\textbf{Problème 2}} \\
   BAF décide d'étendre son marché de savons à l'extérieur de la ville. En se servant de la carte géographique ci-dessous,  la société décide de commencer par la ville ayant pour coordonnées géographiques $(30\degres E,30\degres N)$.  \\
\includegraphics[width=0.86\textwidth]{imagen_1_1605032963}
\item Définis les termes suivants : longitude d'un point, latitude d'un point, équateur, méridien.
\item Reproduis puis complète le tableau suivant : \\
\begin{tabular}{|p{2cm}|p{2.5cm}|p{3cm}|p{3cm}|p{3cm}|}
\hline 
Villes & La Nouvelle-Orléans & Shanghai & Londres & Le Caire \\
\hline
Longitude & & & & \\
\hline
Latitude & & & & \\
\hline
\end{tabular}
\item Déduis-en la ville dans laquelle BAF veut commencer l'extension de son marché de savon. \\ \\
\underline{\textbf{Problème 3}} \\ Les voies $(D1)$, $(D2)$, $(D3)$ et $(D4)$ représentent des droites et les points de ventes sont marqués par les lettres $A, T, J, K, B, H, L$.
\item 
\begin{enumerate}
\item Recopie puis complète par $\in$ ou $\notin$ : $T...[BL)$ ; $A...(D3)$ ; $L...(D4)$ 
\item Donne un autre nom aux droites $(D1)$ et $(D2)$ puis cite deux demi-droites opposés sur le schéma.
\end{enumerate}
\item Identifie deux droites sécantes sur le schéma puis précise leur point d'intersection.
\item Démontre à l'aide d'un organigramme que $(D2) \parallel (D3)$.
\end{enumerate}
\end{document}
